\documentclass[11pt,a4paper]{article}

\usepackage[spanish,es-noquoting]{babel}
\usepackage[utf8]{inputenc}
\usepackage[T1]{fontenc}
\usepackage{geometry}
\usepackage[hyphens,spaces,obeyspaces]{url}
\usepackage{hyperref}
\hypersetup{breaklinks=true} % Permite el salto de línea en los enlaces
\usepackage{xcolor}

\definecolor{raro-jwt-1}{RGB}{71, 61, 16}
\definecolor{raro-jwt-2}{RGB}{130, 69, 66}

\usepackage{adjustbox}
\usepackage{listings}
\lstdefinelanguage{JavaCustom}{
language=Java,
morekeywords={var}
}

\lstdefinelanguage{json}{
language=Java
}

\lstdefinelanguage{properties}{
  sensitive=false,
  morecomment=[l]{\#},
  morecomment=[l]{!},
  morestring=[b]"
}

% Definir el estilo y la sintaxis para YAML
\lstdefinelanguage{yaml}{
  keywords={true, false, null, y, n},
  keywordstyle=\color{blue}\bfseries,
  comment=[l]{\#},
  commentstyle=\color{gray}\ttfamily,
  stringstyle=\color{orange}\ttfamily,
  sensitive=false,
  morestring=[b]',
  morestring=[b]",
  alsoletter={-},
}

\lstset{
    inputencoding=utf8,
    extendedchars=true,
    literate=
        {á}{{\'a}}1
        {é}{{\'e}}1
        {í}{{\'i}}1
        {ó}{{\'o}}1
        {ú}{{\'u}}1
        {Á}{{\'A}}1
        {É}{{\'E}}1
        {Í}{{\'I}}1
        {Ó}{{\'O}}1
        {Ú}{{\'U}}1
        {ñ}{{\~n}}1
        {Ñ}{{\~N}}1,
basicstyle=\ttfamily\small,
keywordstyle=\color{blue}\bfseries,
stringstyle=\color{orange},
commentstyle=\color{gray}\itshape,
numberstyle=\tiny\color{gray},
identifierstyle=\color{black},
%backgroundcolor=\color{gray!10},
frame=none,
breaklines=true,
showstringspaces=false,
tabsize=1,
keepspaces=true,
columns=fullflexible,
numbers=none,
numbersep=8pt,
rulecolor=\color{black},
captionpos=b,
xleftmargin=1.5em,
framexleftmargin=1.5em
}

\usepackage{natbib}
\usepackage{graphicx}
\usepackage{endnotes}
\renewcommand{\notesname}{Notas finales}

\usepackage{tikz}
\usetikzlibrary{arrows.meta, positioning, backgrounds, fit, babel}
\pgfdeclarelayer{background}
\pgfsetlayers{background,main}

\usepackage{seqsplit}
\usepackage{float}

\geometry{margin=2.5cm}

\title{Gestión de \textit{Secretos} (contraseñas, rutas, etc.)\\AWS Secrets Manager}
\author{}
\date{\today}

\begin{document}

%\author{Oscar Dieste} % Your name

\date{\normalsize\today} % Today's date or a custom date

\maketitle

%\begin{abstract}
%\end{abstract}

{
\let\clearpage\relax
\tableofcontents
}

\section{Gestión manual de \textit{secretos}}

Por \textit{secreto}, debemos entender cualquier tipo de información que debe permanecer oculta, ya que su difusión implicaría algún riesgo, normalmente el acceso indebido. Los ejemplos típicos son las contraseñas que permiten acceder a ciertos servicios (como es el caso de las bases de datos), las API KEYs que permiten el acceso a APIs externas, tokens OAuth/M2M (machine-to-machine) utilizados por una aplicación para autenticarse ante ciertos microservicios, las claves privadas, etc.

Excluimos del concepto de secreto las claves de usuario. Desde luego, una clave de usuario debe mantenerse en secreto. Sin embargo, cuando una aplicación requiere una clave de usuario, el usuario puede proporcionarla introduciéndola directamente por teclado. Sin embargo, en el caso de los \textit{secretos}, la situación es distinta porque quien debe poseer ese secreto es algún tipo de aplicación, por ejemplo, el programa que necesita acceder a una base de datos. Esto exige que la clave de acceso esté disponible en algún lugar (no hay usuario que la pueda proporcionar), lo que propicia su acceso/interceptación y uso indebido.

Antes de contar con gestores centralizados (como es el caso de AWS \textit{Secrets Manager}, que introduciremos posteriormente), el almacenamiento de credenciales solía gestionarse mediante tres enfoques, cada uno con importantes vulnerabilidades:

\begin{itemize}
    \item \textbf{\textit{Hardcode} en el código fuente:} Los secretos (API KEYs, contraseñas, tokens, etc.) se incluyen directamente en el código, por ejemplo:

\begin{lstlisting}[language=Java,basicstyle=\ttfamily\footnotesize]
package com.ejemplo.db;

import java.sql.Connection;
import java.sql.DriverManager;
import java.sql.SQLException;

public class DatabaseConnector {

    /*
    * Secretos
    */
    private static final String DB_URL = "jdbc:postgresql://db.ejemplo.com:5432/midatabase";
    private static final String DB_USER = "admin_user";
    private static final String DB_PASSWORD = "Password_1234";

    public static void main(String[] args) {
        try (Connection conn = DriverManager.getConnection(DB_URL, DB_USER, DB_PASSWORD)) {
            ...
        } catch (SQLException e) {
            System.err.println("Error de conexión: " + e.getMessage());
        }
    }
}
\end{lstlisting}
    
    Este procedimiento es inseguro ya que cualquier persona con acceso al código fuente puede acceder al secreto. El secreto también podría recuperarse de una versión ejecutable, ya que los strings aparecen en el segmento de datos del código. 

    \item \textbf{Ficheros de configuración locales (\texttt{.env}, \texttt{.json}, \texttt{.yaml}):} Los secretos se guardan en ficheros (normalmente de texto) dentro del servidor o entorno local de cada desarrollador, como podría ser el siguiente \texttt{config.properties}:

\begin{lstlisting}[language=Properties,basicstyle=\ttfamily\footnotesize]
db.url=jdbc:postgresql://db.ejemplo.com:5432/midatabase
db.user=admin_user
db.password=Password_1234
\end{lstlisting}

Su uso sería, por ejemplo, el siguiente:

\begin{lstlisting}[language=Java, basicstyle=\ttfamily\footnotesize]
package com.ejemplo.db;

import java.io.IOException;
import java.io.InputStream;
import java.sql.Connection;
import java.sql.DriverManager;
import java.sql.SQLException;
import java.util.Properties;

public class DatabaseConnector {

    public static void main(String[] args) {
        Properties properties = new Properties();

        // Cargar el archivo de propiedades desde el classpath
        try (InputStream input = DatabaseConnector.class.getClassLoader().getResourceAsStream("config.properties")) {
            properties.load(input);
        } catch (IOException e) {
            System.err.println(e.getMessage());
            return;
        }

        // Extraer las credenciales y conectar
        String dbUrl = properties.getProperty("db.url");
        String dbUser = properties.getProperty("db.user");
        String dbPassword = properties.getProperty("db.password");

        try (Connection conn = DriverManager.getConnection(dbUrl, dbUser, dbPassword)) {
            ...
        } catch (SQLException e) {
            System.err.println(e.getMessage());
        }
    }
}
\end{lstlisting}

    Aunque no se suban al repositorio (por ejemplo, usando \texttt{.gitignore}), estos ficheros quedan dispersos en los equipos de desarrollo sin, probablemente, un control estricto de acceso. Además, los ficheros de configuración deben subirse a los servidores%
    \footnote{Subir los ficheros a los servidores introduce un problema adicional: 1) O bien dichos archivos se transfieren manualmente, por ejemplo, usando \texttt{sftp}, o bien 2) habría que utilizar un pipeline CI/CD local en alguna máquina de desarrollo.     
    Una alternativa intermedia sería confiar en pipelines CI/CD públicos con seguridad contrastada, tales como GitHub, Bitbucket o CodeDeploy, o un Jenkins institucional bien protegido. Sin embargo, en todos estos casos anteriores, volvemos a tener el problema de que los secretos estarían disponibles en los repositorios.}, %
    por lo que cualquiera que tenga acceso a los servidores podría obtener los secretos.

    \item \textbf{Variables de entorno:} Definir las variables directamente en la configuración del servidor web o del contenedor (que es algo muy frecuente), por ejemplo:

\begin{lstlisting}[basicstyle=\ttfamily\footnotesize]

# Imagen base oficial de AWS para Java (Amazon Corretto 21)
FROM amazoncorretto:21-alpine

WORKDIR /app

ENV DB_URL="jdbc:postgresql://db.ejemplo.com:5432/midatabase" 
ENV DB_USER="admin_user"
ENV DB_PASSWORD="Password_1234"

COPY target/*.jar app.jar
EXPOSE 8080
ENTRYPOINT ["java", "-jar", "app.jar"]
\end{lstlisting}

\vspace{1em}

las cuales se podrían utilizar de la forma siguiente:

\vspace{1em}

\begin{lstlisting}[language=Java,basicstyle=\ttfamily\footnotesize]
package com.ejemplo.db;

import java.sql.Connection;
import java.sql.DriverManager;
import java.sql.SQLException;

public class DatabaseConnector {

    /*
    * Secretos
    */
    private static final String DB_URL = System.getenv("DB_URL");
    private static final String DB_USER = System.getenv("DB_USER");
    private static final String DB_PASSWORD = System.getenv("DB_PASSWORD");

    public static void main(String[] args) {
        try (Connection conn = DriverManager.getConnection(DB_URL, DB_USER, DB_PASSWORD)) {
            ...
        } catch (SQLException e) {
            System.err.println("Error de conexión: " + e.getMessage());
        }
    }
}
\end{lstlisting}

Los problemas de las variables de entorno son muy parecidas a los problemas de los archivos de configuración.

\end{itemize}

\section{AWS Secrets Manager}

\textbf{AWS Secrets Manager} es un \textit{vault} (bóveda de seguridad) administrado por AWS y diseñado para almacenar, recuperar y auditar secretos a lo largo de su ciclo de vida, evitando la necesidad de mantenerlos estáticos en disco o en memoria. Existen otras soluciones en otras plataformas como GCP Secret Manager, Azure Key Vault, etc.

Lo más importante para nuestros propósitos es que el control de acceso a los secretos se gestiona mediante políticas de permisos nativas de AWS Identity and Access Management (IAM), limitando qué usuario o recurso (como una instancia EC2 o una función Lambda) puede descifrar cada secreto específico. Lo veremos en la sección~\ref{sec:uso}.

En el resto del documento usaremos como ejemplo el \textit{backend} Spring Boot de Event Hub. Su base de datos es H2, una base de datos relacional escrita en Java que se ejecuta dentro del propio proceso de la aplicación (embebida) y guarda los datos en memoria. Al ser embebida, H2 no tiene servidor, ni \textit{host}, ni puerto: el secreto solo necesita el usuario, la contraseña y el nombre de la base de datos.

H2 tiene una propiedad que simplifica el ejemplo: la base de datos en memoria se crea en la primera conexión, con el usuario y la contraseña que recibe en ese momento. Así, basta con que Spring Boot lea el secreto y se conecte; no hay que dar de alta el usuario en el gestor de base de datos. Con una base de datos con servidor (PostgreSQL, MySQL, etc.), después de crear el secreto habría que hacer un \texttt{CREATE USER} (o un \texttt{ALTER USER} si estamos cambiando la contraseña) para que el gestor conozca esa credencial.

\subsection{Creación de secretos mediante la consola web}

AWS Secrets Manager puede usarse directamente a través de la consola web de AWS. Sólo es necesario acceder al servicio y definir los secretos. La figura~\ref{fig:pantalla-secrets-manager} muestra el interfaz de usuario.

\begin{figure}
    \centering
    \includegraphics[width=\linewidth]{figs/pantalla-secrets-manager.png}
    \caption{Pantalla de Secrets manager que permite la creación de secretos}
    \label{fig:pantalla-secrets-manager}
\end{figure}

El interfaz de usuario sugiere que hay distintos tipos de secretos. Esto no es correcto. Un secreto es simplemente un objeto json formado por pares clave-valor. La pantalla de credenciales de base de datos funciona únicamente como una plantilla que genera un objeto json con campos predefinidos como \texttt{engine}, \texttt{host}, \texttt{username}, \texttt{password}, \texttt{dbname} y \texttt{port}. No es necesario respetar esta plantilla: el secreto puede contener las claves que necesite la aplicación.

Utilizar la consola web es una muy buena opción en muchos casos. Un administrador podría definir los secretos, y estos nunca estarían disponibles en texto claro. Los secretos se podrían utilizar tal y como se indica en la sección~\ref{sec:uso}. Esto funciona bastante bien para bases de datos, API KEYs, y secretos que tienen permanencia en el tiempo.

\subsection{Generación de secretos mediante el CLI}

Utilizando el CLI, los secretos podrían aparecer explícitamente en los scripts, por ejemplo:

\begin{lstlisting}[language=bash]
# Variables de configuracion (ejemplo generico)
DB_USER="admin"
DB_PASS="Un4 Contr4senn4 d|f|c|l;"
...
\end{lstlisting}

\vspace{1em}

Pero esto supondría, de nuevo, que los secretos se almacenan en texto claro. En su lugar, hay que generar los secretos \textbf{automáticamente}. La mejor opción es usar \texttt{aws secretsmanager get-random-password}. El script \texttt{aws-scripts/secrets.sh} de Event Hub lo hace así:

\begin{lstlisting}[language=bash]
# Credenciales de H2 para Spring Boot. H2 es embebida: no hay host ni puerto.
SECRET_NAME="prod/h2/admin"

# En H2 el usuario puede ser cualquiera.
DB_USER="admin"
DB_NAME="eventhub"

# ExcludeCharacters: /, ", @, ' y \ suelen romper JSON, JDBC, SQL o el shell.
DB_PASS=$(aws secretsmanager get-random-password \
    --password-length 32 \
    --exclude-characters "/\"@'\\" \
    --query RandomPassword \
    --output text)
if [ -z "$DB_PASS" ] || [ "$DB_PASS" = "None" ]; then
    echo "Error: No se pudo generar la password."
    exit 1
fi

SECRET_JSON=$(jq -c -n \
    --arg user "$DB_USER" \
    --arg pass "$DB_PASS" \
    --arg dbname "$DB_NAME" \
    '{username: $user, password: $pass, dbname: $dbname}')
unset DB_PASS
\end{lstlisting}

\vspace{1em}

\texttt{SECRET\_NAME} constituye el identificador único utilizado por AWS Secrets Manager para registrar, organizar y recuperar un secreto dentro de una cuenta y región determinadas. Se recomienda emplear la convención de ruta *nix (\texttt{foo/bar}), ya que es sencilla de entender y permite restringir permisos en AWS IAM mediante comodines (como \texttt{prod/h2/*}). Hay que tener en cuenta que Secrets Manager trata el nombre como una cadena literal: \texttt{prod/h2/admin} y \texttt{/prod/h2/admin} son dos secretos distintos.

\texttt{\textquotedbl/\textbackslash\textquotedbl @\textquotesingle\textbackslash\textbackslash\textquotedbl} es como se escribe en el shell la lista de caracteres /, ", @, ' y \textbackslash. Estos caracteres se excluyen porque suelen romper JSON, URLs JDBC, SQL o el mismo shell si la password se embebe en un string.

\texttt{get-random-password} genera la contraseña en AWS y la devuelve al script; no se crea ningún recurso. Si la llamada falla, el script se detiene antes de crear un secreto con la contraseña vacía.

Finalmente, \texttt{SECRET\_JSON} es el objeto json con pares clave-valor del que hemos hablado anteriormente. \texttt{jq} lo construye a partir de las variables, lo que evita problemas de comillas si la contraseña contiene caracteres especiales; \texttt{-n} indica que no hay entrada y \texttt{-c} genera el JSON en una sola línea. En cuanto la contraseña está dentro del JSON, \texttt{unset DB\_PASS} la elimina de la variable. Aunque hemos creado \texttt{SECRET\_JSON}, todavía no lo hemos subido a Secrets Manager. Esto es, el JSON no ha salido de la máquina local.

Hemos utilizado como ejemplo las credenciales de acceso a una base de datos. Sin embargo, el mismo procedimiento serviría para cualquier tipo de secreto.

\subsection{Creación de los secretos usando el CLI}\label{sec:creacion-secreto}

Finalmente, nos queda crear, o subir si se quiere, el secreto a Secrets Manager:

\begin{lstlisting}[language=bash]
# --secret-string: JSON con username, password y dbname. No se imprime.
SECRET_ARN=$(aws secretsmanager create-secret \
    --name "$SECRET_NAME" \
    --description "Credenciales de administración ($DB_USER) para la base de datos '$DB_NAME'." \
    --secret-string "$SECRET_JSON" \
    --query ARN \
    --output text)
unset SECRET_JSON
\end{lstlisting}

\vspace{1em}

\texttt{--query ARN} hace que el CLI devuelva solo el identificador del secreto (\textit{Amazon Resource Name}) y no la respuesta completa. El script muestra el nombre y el ARN, nunca el contenido.

El borrado es la operación inversa (\texttt{secrets.sh delete}):

\begin{lstlisting}[language=bash]
# --force-delete-without-recovery: borrado inmediato, sin ventana de recuperación
aws secretsmanager delete-secret \
    --secret-id "$SECRET_NAME" \
    --force-delete-without-recovery
\end{lstlisting}

\vspace{1em}

Por defecto, Secrets Manager no borra un secreto de inmediato: lo programa para borrarlo tras una ventana de recuperación de entre 7 y 30 días. Mientras tanto, el nombre sigue ocupado y un nuevo \texttt{create-secret} con el mismo nombre fallaría. En un laboratorio que se crea y se destruye a menudo, \texttt{--force-delete-without-recovery} evita ese problema.

En Event Hub, \texttt{make deploy} ejecuta \texttt{secrets.sh create} después de crear la instancia EC2 y antes de desplegar el \textit{backend}, y \texttt{make delete} ejecuta \texttt{secrets.sh delete}.

\subsection{Uso de los secretos}\label{sec:uso}

Independientemente del modo en que los secretos han sido creados, los pasos para ser usados son los mismos en ambos casos.

\subsubsection{Permisos y credenciales}

En primer lugar, quien lee el secreto necesita credenciales de AWS con permiso para ejecutar \texttt{secretsmanager:GetSecretValue} sobre él. La asignación de permisos es un tema complejo para exponer aquí, pero está explicado en detalle en la \href{https://www.overleaf.com/read/tfyjkhsjhztf#20828d}{Guía de IAM}. Además, el ejemplo que se plantea en dicha guía es, precisamente, la asignación de permiso de lectura de secretos a una función Lambda y a una máquina EC2.

El SDK de AWS (y, por tanto, Spring Boot) obtiene las credenciales recorriendo una cadena de proveedores, en este orden aproximado:

\begin{enumerate}
    \item Variables de entorno (\texttt{AWS\_ACCESS\_KEY\_ID}, \texttt{AWS\_SECRET\_ACCESS\_KEY}, \texttt{AWS\_SESSION\_TOKEN}).
    \item El fichero \texttt{\textasciitilde/.aws/credentials}, el mismo que usa el CLI.
    \item Las credenciales del contenedor (ECS) o del rol asignado a la instancia EC2, que el SDK obtiene del servicio de metadatos de la instancia.
\end{enumerate}

En consecuencia, la aplicación puede leer el secreto tanto en la máquina de desarrollo (con las credenciales del CLI) como en una instancia EC2. La diferencia está en que, en la EC2, las credenciales no se guardan en disco: las proporciona el rol asignado a la instancia, y son temporales y se renuevan automáticamente.

En AWS Academy no es posible crear roles. El laboratorio ya incluye el rol \texttt{LabRole} y su perfil de instancia \texttt{LabInstanceProfile}, que se asocian a la instancia al crearla. En Event Hub lo hace \texttt{aws-scripts/ec2.sh}:

\begin{lstlisting}[language=bash]
# --iam-instance-profile: rol LabRole de AWS Academy. Spring Boot lo usa
# para leer el secreto en Secrets Manager.
INSTANCE_ID=$(aws ec2 run-instances \
    --image-id "$AMI_ID" \
    --count 1 \
    --instance-type "$INSTANCE_TYPE" \
    --key-name "$KEY_NAME" \
    --iam-instance-profile Name="$INSTANCE_PROFILE" \
    --security-group-ids "$SG_ID" \
    --query "Instances[0].InstanceId" \
    --output text)
\end{lstlisting}

\vspace{1em}

\texttt{INSTANCE\_PROFILE} vale \texttt{LabInstanceProfile}. Sin esta opción, la instancia no tiene credenciales y la aplicación no puede leer el secreto.

\subsubsection{Lectura con el SDK de AWS}

Una vez asignados los permisos, el recurso que haya sido autorizado puede ejecutar la consulta sin necesidad de configurar claves. Por ejemplo, se podría consultar el secreto mediante el siguiente código Java (se muestra como un servicio de Spring Boot, pero no sería necesario hacerlo así; funcionaría exactamente igual, por ejemplo, dentro de una función \texttt{main()}):

\begin{lstlisting}[language=Java]
package com.example.demo.service;

import org.springframework.stereotype.Service;
import software.amazon.awssdk.core.exception.SdkClientException;
import software.amazon.awssdk.services.secretsmanager.SecretsManagerClient;
import software.amazon.awssdk.services.secretsmanager.model.GetSecretValueRequest;
import software.amazon.awssdk.services.secretsmanager.model.SecretsManagerException;

@Service
public class SecretService {

    private final SecretsManagerClient secretsClient;

    public SecretService() {
        // Sin credenciales ni región explícitas, el SDK recorre la cadena de
        // proveedores: variables de entorno, ~/.aws, rol de la EC2, etc.
        this.secretsClient = SecretsManagerClient.create();
    }

    // SecretsManagerException: la ejecución falla por falta de permisos IAM
    // o inexistencia del secreto
    // SdkClientException: la ejecución falla por problemas locales, normalmente
    // falta de credenciales o de conectividad a la red
    public String obtenerCredencialesDb() throws SecretsManagerException, SdkClientException {
        GetSecretValueRequest request = GetSecretValueRequest.builder()
                .secretId("prod/h2/admin")
                .build();

        return secretsClient.getSecretValue(request).secretString();
    }
}
\end{lstlisting}

\texttt{SecretsManagerClient} y \texttt{GetSecretValueRequest} son clases del SDK de AWS que permiten acceder a los secretos. \texttt{secretString()} devuelve el JSON del secreto como texto; la aplicación tendría que interpretarlo para extraer cada campo. Si no hay credenciales disponibles o no tienen permiso sobre el secreto, el código fallará con alguna de las excepciones anteriores.

\subsubsection{Lectura con Spring Cloud AWS}

Event Hub no usa el SDK directamente. Cuando usamos Spring Boot, la alternativa más cómoda es la librería \textit{starter} \texttt{spring-cloud-aws-starter-secrets-manager}, que consulta los secretos durante la inicialización y los añade al entorno de Spring, de modo que pueden usarse como cualquier otra propiedad. La aproximación es:

\begin{itemize}
    \item Añadir el BOM y la dependencia Maven (o Gradle). En el \texttt{pom.xml} de Event Hub:

\begin{lstlisting}[language=XML]
<properties>
    <spring-cloud-aws.version>4.1.0</spring-cloud-aws.version>
</properties>

<dependencyManagement>
    <dependencies>
        <dependency>
            <groupId>io.awspring.cloud</groupId>
            <artifactId>spring-cloud-aws-dependencies</artifactId>
            <version>${spring-cloud-aws.version}</version>
            <type>pom</type>
            <scope>import</scope>
        </dependency>
    </dependencies>
</dependencyManagement>

<dependency>
    <groupId>io.awspring.cloud</groupId>
    <artifactId>spring-cloud-aws-starter-secrets-manager</artifactId>
</dependency>
\end{lstlisting}

El BOM \texttt{spring-cloud-aws-dependencies} fija la versión del \textit{starter}, por lo que la dependencia no la indica. La versión 4.1.0 es la compatible con la versión de Spring Boot del proyecto (4.1).

\item Indicar la importación del secreto utilizando la opción \texttt{aws-secretsmanager:}. En el \texttt{application.properties} de Event Hub:

\begin{lstlisting}[language=properties]
spring.config.import=classpath:cognito.properties,aws-secretsmanager:prod/h2/admin
\end{lstlisting}

\texttt{spring.config.import} admite varias fuentes separadas por comas. \texttt{classpath:cognito.properties} es un fichero que genera otro script del laboratorio; \texttt{aws-secretsmanager:prod/h2/admin} es el secreto. El nombre tiene que coincidir exactamente con el que se usó en \texttt{create-secret}, sin barra inicial. Si el secreto no existe o no hay permisos, la aplicación no arranca.

\item Usar los valores. Cada clave de primer nivel del JSON del secreto (\texttt{username}, \texttt{password} y \texttt{dbname}) se convierte en una propiedad de Spring con el mismo nombre. Por ejemplo, con la anotación \texttt{@Value}:

\begin{lstlisting}[language=Java]
public class DatabaseHealthService {

    @Value("${username}")
    private String dbUser;

    public void logConnectionUser() {
        // El valor proviene directamente del secreto prod/h2/admin
        System.out.println("Usuario de BD configurado: " + dbUser);
    }
}
\end{lstlisting}

\end{itemize}

\subsubsection{Configuración de la base de datos}

Si el secreto utilizase las mismas claves estándar del framework (por ejemplo, \texttt{spring.datasource.username}), Spring Boot podría configurar la conexión a la base de datos de forma automática. En Spring Boot, las claves estándar son:

\begin{itemize}
\item \texttt{spring.datasource.url}: Cadena de conexión JDBC.
\item \texttt{spring.datasource.username}: Nombre del usuario.
\item \texttt{spring.datasource.password}: Contraseña.
\item \texttt{spring.datasource.driver-class-name}: Driver JDBC (\textit{opcional}, suele inferirse de la URL).
\end{itemize}

Como las claves estándar y las claves del secreto no coinciden, tenemos que aplicar una de estas opciones: 1) usar las claves de Spring Boot en Secrets Manager (opción perfectamente viable, pero que ata el secreto a un \textit{framework} concreto) o 2) mapear las claves del secreto a claves de Spring Boot en el fichero \texttt{application.properties}. Event Hub usa la segunda:

\begin{lstlisting}[language=properties]
spring.datasource.driverClassName=org.h2.Driver
spring.datasource.username=${username}
spring.datasource.password=${password}
spring.datasource.url=jdbc:h2:mem:${dbname};DB_CLOSE_DELAY=-1;DB_CLOSE_ON_EXIT=FALSE
\end{lstlisting}

\begin{itemize}
    \item \texttt{\$\{username\}}, \texttt{\$\{password\}} y \texttt{\$\{dbname\}} se resuelven con las claves del secreto.
    \item \texttt{jdbc:h2:mem:eventhub} indica una base de datos H2 en memoria llamada \texttt{eventhub}. Los datos se pierden al detener la aplicación.
    \item \texttt{DB\_CLOSE\_DELAY=-1} mantiene la base de datos mientras la aplicación esté en marcha, aunque no quede ninguna conexión abierta.
    \item \texttt{DB\_CLOSE\_ON\_EXIT=FALSE} deja que sea Spring, y no H2, quien cierre la base de datos al terminar.
\end{itemize}

En la primera conexión, H2 crea la base de datos \texttt{eventhub} con el usuario y la contraseña del secreto. Una vez inicializada la base de datos, podría usarse normalmente utilizando JPA, o usando el bean \texttt{JdbcTemplate} en el caso de que se deseen realizar invocaciones directas en SQL.

Por último, los tests no deberían depender de AWS. Event Hub tiene un \texttt{application.properties} propio en \texttt{src/test/resources} que no importa el secreto y usa H2 con credenciales fijas, de modo que \texttt{mvn test} funciona sin credenciales:

\begin{lstlisting}[language=properties]
# Tests: H2 en memoria con credenciales fijas. No se importa Secrets Manager.
spring.config.import=classpath:cognito.properties
spring.datasource.url=jdbc:h2:mem:eventhubdb;DB_CLOSE_DELAY=-1;DB_CLOSE_ON_EXIT=FALSE
spring.datasource.driverClassName=org.h2.Driver
spring.datasource.username=sa
spring.datasource.password=
\end{lstlisting}

%\theendnotes

\end{document}
